\documentclass[11pt]{article}
\usepackage[margin=1in]{geometry}
\usepackage{booktabs}
\usepackage{longtable}
\usepackage{array}
\usepackage{fontspec}
\usepackage{hyperref}
\usepackage{enumitem}
\usepackage{setspace}
\setmainfont{Noto Sans CJK KR}
\setlength{\parskip}{0.45em}
\setlength{\parindent}{0pt}
\tolerance=1000
\emergencystretch=3em
\hfuzz=2pt
\title{Meta-CoT 실험 중간 보고서}
\author{Seungpil Lee}
\date{2026-04-01}

\begin{document}
\maketitle
\onehalfspacing

\section*{요약}
현재까지의 증거는 Meta-CoT가 메타 텍스트를 생성하는 능력 자체는 확보했지만, 우리가 원한 메타인지 제어 정책 전체를 아직 안정적으로 학습한 것은 아니라는 점을 보여준다. 특히 hard problem에서 wrong-answer confidence를 낮추는 방향의 보상은 일부 효과가 있었지만, 그 자체만으로 AIME 정확도가 안정적으로 좋아지지는 않았다. 따라서 현재의 핵심 과제는 메타 표현을 더 늘리는 것이 아니라, confidence가 실제 행동 전환으로 이어지도록 만드는 것이다.

현재 계획의 중심 의도는 두 가지다. 첫째, 막힘이나 모순, anomaly를 감지했을 때 confidence를 낮추고 실제로 다른 풀이 방법으로 redirect해야 한다. 둘째, 높은 confidence를 가질 때에도 final answer 전에 독립적 verify를 수행해야 한다. 지금의 실험 구성은 이 두 의도를 기준으로 V2/V3, E-series, behavior-first SFT, 차기 E10 비교로 분해되어 있으며, 이 구조는 연구 의도와 대체로 잘 정렬되어 있다.

\section*{핵심 의도}
이 프로젝트의 목적은 더 많은 meta text를 생성하는 것이 아니다. 목적은 confidence를 내부 제어 신호로 사용하여 reasoning trajectory를 바꾸는 것이다. 따라서 성공 기준은 ``메타를 말했다''가 아니라 ``적절한 조건에서 행동이 바뀌었다''가 되어야 한다.

우리가 원하는 첫 번째 행동은 failure-management다. contradiction, failed substitution, unsupported assumption, unit mismatch, 또는 confidence drop이 나타나면 현재 경로를 밀어붙이지 않고 왜 안 풀리는지 짧게 진단한 뒤 다른 방법으로 redirect해야 한다. 두 번째 행동은 overconfidence-management다. 모델이 자신 있어 보일수록 final answer 전에 independent verification을 수행해야 한다.

\section*{정량 결과}
\begin{center}
\begin{tabular}{lr}
\toprule
모델 & 1,030 문제 정확도 \\
\midrule
Base SFT & 71.70\% \\
V2 SFT & 72.72\% \\
V3 SFT & 72.00\% \\
E5 & 72.04\% \\
E7 prev & 69.90\% \\
E7 current & 70.68\% \\
E8 & 67.20\% \\
V2 rich & 71.36\% \\
\bottomrule
\end{tabular}
\end{center}

이 표는 메타 구조가 수학 성능과 반드시 충돌하지는 않는다는 점을 보여준다. 다만 이것만으로는 메타 제어가 학습됐다고 말할 수 없다. 일부 모델은 base 수준을 유지하거나 약간 넘지만, 더 많은 meta intervention이 항상 더 나은 제어를 뜻하지는 않는다.

\begin{center}
\begin{tabular}{lrrr}
\toprule
모델 & AIME & 오답 평균 confidence & 오답 평균 meta block 수 \\
\midrule
Base SFT & 3/30 & N/A & 0.00 \\
V2 SFT & 4/30 & 0.631 & 2.58 \\
V3 SFT & 2/30 & 0.287 & 2.32 \\
E5 & 2/30 & 0.321 & 5.36 \\
E7 prev & 3/30 & 0.315 & 5.15 \\
E7 current & 1/30 & 0.263 & 4.34 \\
E8 & 2/30 & 0.238 & 5.04 \\
\bottomrule
\end{tabular}
\end{center}

이 표는 calibration 축의 부분적 성공과 한계를 동시에 보여준다. RL 계열 모델은 hard wrong answer에서 V2보다 훨씬 덜 과신하는 경향을 보인다. 그러나 낮아진 confidence와 늘어난 meta block이 아직 진짜 diagnosis-driven control로 이어졌다고 보기는 어렵다.

\section*{정성 해석}
이전 정성 분석에서 관찰된 핵심 패턴은 다음과 같다. V2는 중간 정도의 confidence를 유지한 채 비교적 관습적인 풀이를 지속하는 경우가 많다. 반면 E7과 E8은 더 자주 스스로를 끊고 confidence를 낮춘다. 이 자체는 실제 행동 변화의 초기 신호로 볼 수 있다.

문제는 그 다음이다. 현재의 개입은 종종 local correction이나 self-talk 수준에 머문다. 즉, ``조심하자'' 또는 ``다른 방법을 보자''라고 말하지만, 왜 현재 경로가 실패했는지 명시적으로 진단하지 않거나, 필요한 subgoal을 다시 세우지 않거나, 실제로 다른 풀이 전략으로 이동하지 않는 경우가 많다. 따라서 현재 결론은 ``아이디어가 틀렸다''가 아니라 ``calibration과 interruption은 일부 학습됐지만, diagnosis와 redirect는 아직 약하다''가 더 정확하다.

\section*{현재 실험 구성의 정렬성}
현재 실험 구성은 각 단계의 역할이 분리되어 있다는 점에서 연구 의도와 잘 맞는다. V2와 V3는 explicit meta 표현과 confidence 표기가 가능한지 확인하는 표현 단계다. E3와 E8은 confidence가 실제 오류 위험과 더 잘 정렬되는지 보는 calibration 단계다. Behavior-first SFT는 verify와 redirect를 직접 행동으로 가르치는 제어 단계다. 차기 E10은 calibration 축을 버리지 않고 behavior reward를 위에 쌓아, 어떤 보상 축이 실제 행동 변화를 만드는지 분해해 볼 수 있는 단계가 되어야 한다.

이 정렬성은 중요하다. 지금 curriculum이나 retrieval을 먼저 붙이면, 이후 향상이 self-diagnosis에서 온 것인지 외부 보조에서 온 것인지 분리하기 어렵다. 따라서 현재처럼 control policy를 먼저 안정화하고 그 다음 curriculum/RAG로 넘어가는 순서가 맞다.

\section*{파일럿 데이터의 의미}
Behavior-first pilot은 방향 자체는 옳다는 점을 보여줬다. 이제 프로젝트는 ``straight'', ``verify'', ``redirect''라는 구체적 행동 단위로 supervision을 설계할 수 있다. 그러나 현재 pilot은 1,800개 요청 중 770개만 유효했고, 그중 552개가 redirect, 216개가 verify, 2개만 straight였다. 이 상태로는 calibrated controller가 아니라 over-intervention model을 만들 위험이 크다.

따라서 다음 단계는 단순한 스케일업이 아니다. 우선 natural meta intervention 형식을 유지하면서도 straight, verify, redirect가 모두 살아남는 control-v4 pilot을 다시 만들어야 한다. 그 후에 base\_sft에서 출발하는 behavior\_verify\_sft\_v2, behavior\_redirect\_sft\_v2, behavior\_all\_sft\_v2를 비교하고, 마지막으로 E3, E8, E10의 RL 비교로 넘어가야 한다.

\section*{현재 결론}
현재 계획 문서에 들어 있는 중심 의도는 올바르다. low confidence에서의 redirect와 high confidence에서의 verify를 분리해 정의한 점, confidence를 scalar가 아니라 control signal로 해석한 점, curriculum을 그 이후 단계로 게이트한 점이 모두 같은 방향을 가리킨다. 지금 필요한 것은 방향 전환이 아니라 구현 정교화다. 데이터는 더 자연스럽고 균형 있게 만들어야 하고, reward는 calibration 축과 behavior 축을 분해해서 비교해야 하며, 분석은 meta text가 아니라 행동 변화를 중심으로 설계해야 한다.

\end{document}
